% GENERATED FILE. Edit canonical lessons, then run npm run generate:books.
% canonical-source-hash: fnv1a64:2172abe148630283
% canonical-lessons: SA-C203-modate, SA-C203-vandate, SA-C203-yacate, SA-C203-bhasate, SA-C203-sobhate

\chapter{Rejoices, Salutes, Begs, Speaks, Shines}
\label{ch:sa-rejoices-salutes-begs-speaks-shines}
\hlchaptermodality{203}

\begin{chapteropening}
\textbf{By the end of this chapter:} \emph{I can say rejoices, salutes, begs, speaks and shines.}

Complete the last lesson of chapter 203: I can say rejoices, salutes, begs, speaks and shines.
\end{chapteropening}

\section[\texorpdfstring{\sk{मोदते} (modate)}{मोदते (modate)}]{\sk{मोदते} (modate) — rejoices}

\label{lesson:SA-C203-modate}



\begin{quote}
Before the new one: say the Sanskrit for the west, then the Sanskrit for the north.
\end{quote}

\subsection*{You'll want to know: \sk{मोदते}}
\textbf{\sk{मोदते}} — \emph{modate} — \textquotedblleft{}rejoices\textquotedblright{}.

\begin{cousinweb}
\emph{modate} is the present of the root √mud, \textquotedblleft{}to be merry, glad, rejoice\textquotedblright{} (MW). The same root gives \emph{mud}, \textquotedblleft{}joy\textquotedblright{}; MW gives it no older source, so its origin is uncertain.
\end{cousinweb}

\begin{grammarlens}[title={What you've built}]
One more word for this chapter.
\end{grammarlens}

\subsection*{Your turn}
\begin{itemize}
\raggedright
  \item \emph{Say it:} \emph{modate}
  \item \emph{Say it:} \emph{modate}, once more
  \item \emph{Say it:} say \emph{modate}
  \item \emph{Recall:} say the Sanskrit for the west, then the Sanskrit for the north, then say \emph{modate} again
\end{itemize}

\begin{tcolorbox}[breakable,colback=teal!4,colframe=teal!35!black,title={Before you move on}]
What does \sk{मोदते} mean? (\textquotedblleft{}Rejoices\textquotedblright{}.) Say it once more.
\end{tcolorbox}
\section[\texorpdfstring{\sk{वन्दते} (vandate)}{वन्दते (vandate)}]{\sk{वन्दते} (vandate) — salutes, greets with respect}

\label{lesson:SA-C203-vandate}



\begin{quote}
Before the new one: say the Sanskrit for the north, then the Sanskrit for rejoices.
\end{quote}

\subsection*{You'll want to know: \sk{वन्दते}}
\textbf{\sk{वन्दते}} — \emph{vandate} — \textquotedblleft{}salutes, greets with respect\textquotedblright{}.

\begin{cousinweb}
\emph{vandate} is the present of the root √vand, \textquotedblleft{}to praise, salute respectfully\textquotedblright{} (MW). In the Rigveda it means \textquotedblleft{}praises\textquotedblright{}, later \textquotedblleft{}greets respectfully\textquotedblright{} (CDIAL 11270), and its noun \emph{vandana}, \textquotedblleft{}reverence\textquotedblright{}, is borrowed in Malayalam.
\end{cousinweb}

\begin{grammarlens}[title={What you've built}]
One more word for this chapter.
\end{grammarlens}

\subsection*{Your turn}
\begin{itemize}
\raggedright
  \item \emph{Say it:} \emph{vandate}
  \item \emph{Say it:} \emph{vandate}, once more
  \item \emph{Say it:} say \emph{vandate}
  \item \emph{Recall:} say the Sanskrit for the north, then the Sanskrit for rejoices, then say \emph{vandate} again
\end{itemize}

\begin{tcolorbox}[breakable,colback=teal!4,colframe=teal!35!black,title={Before you move on}]
What does \sk{वन्दते} mean? (\textquotedblleft{}Salutes, greets with respect\textquotedblright{}.) Say it once more.
\end{tcolorbox}
\section[\texorpdfstring{\sk{याचते} (yācate)}{याचते (yācate)}]{\sk{याचते} (yācate) — begs, asks for}

\label{lesson:SA-C203-yacate}



\begin{quote}
Before the new one: say the Sanskrit for rejoices, then the Sanskrit for salutes.
\end{quote}

\subsection*{You'll want to know: \sk{याचते}}
\textbf{\sk{याचते}} — \emph{yācate} — \textquotedblleft{}begs, asks for\textquotedblright{}.

\begin{cousinweb}
\emph{yācate} is the middle present of the root √yāc, \textquotedblleft{}to ask, beg\textquotedblright{} (MW): the middle voice means asking for yourself. It is already in the Rigveda (CDIAL 10447), and \emph{yācaka}, \textquotedblleft{}beggar\textquotedblright{}, is built on it.
\end{cousinweb}

\begin{grammarlens}[title={What you've built}]
One more word for this chapter.
\end{grammarlens}

\subsection*{Your turn}
\begin{itemize}
\raggedright
  \item \emph{Say it:} \emph{yācate}
  \item \emph{Say it:} \emph{yācate}, once more
  \item \emph{Say it:} say \emph{yācate}
  \item \emph{Recall:} say the Sanskrit for rejoices, then the Sanskrit for salutes, then say \emph{yācate} again
\end{itemize}

\begin{tcolorbox}[breakable,colback=teal!4,colframe=teal!35!black,title={Before you move on}]
What does \sk{याचते} mean? (\textquotedblleft{}Begs, asks for\textquotedblright{}.) Say it once more.
\end{tcolorbox}
\section[\texorpdfstring{\sk{भाषते} (bhāṣate)}{भाषते (bhāṣate)}]{\sk{भाषते} (bhāṣate) — speaks, talks}

\label{lesson:SA-C203-bhasate}



\begin{quote}
Before the new one: say the Sanskrit for salutes, then the Sanskrit for begs.
\end{quote}

\subsection*{You'll want to know: \sk{भाषते}}
\textbf{\sk{भाषते}} — \emph{bhāṣate} — \textquotedblleft{}speaks, talks\textquotedblright{}.

\begin{cousinweb}
\emph{bhāṣate} is the present of the root √bhāṣ, \textquotedblleft{}to speak, say\textquotedblright{} (MW), first found in the Taittirīya Brāhmaṇa (CDIAL 9478). The nouns \emph{bhāṣā}, \textquotedblleft{}language\textquotedblright{}, and \emph{bhāṣaṇa}, \textquotedblleft{}speech\textquotedblright{}, are built on the same root.
\end{cousinweb}

\begin{grammarlens}[title={What you've built}]
One more word for this chapter.
\end{grammarlens}

\subsection*{Your turn}
\begin{itemize}
\raggedright
  \item \emph{Say it:} \emph{bhāṣate}
  \item \emph{Say it:} \emph{bhāṣate}, once more
  \item \emph{Say it:} say \emph{bhāṣate}
  \item \emph{Recall:} say the Sanskrit for salutes, then the Sanskrit for begs, then say \emph{bhāṣate} again
\end{itemize}

\begin{tcolorbox}[breakable,colback=teal!4,colframe=teal!35!black,title={Before you move on}]
What does \sk{भाषते} mean? (\textquotedblleft{}Speaks, talks\textquotedblright{}.) Say it once more.
\end{tcolorbox}
\section[\texorpdfstring{\sk{शोभते} (śobhate)}{शोभते (śobhate)}]{\sk{शोभते} (śobhate) — shines, looks fine}

\label{lesson:SA-C203-sobhate}



\begin{quote}
Before the new one: say the Sanskrit for begs, then the Sanskrit for speaks.
\end{quote}

\subsection*{You'll want to know: \sk{शोभते}}
\textbf{\sk{शोभते}} — \emph{śobhate} — \textquotedblleft{}shines, looks fine\textquotedblright{}.

\begin{cousinweb}
\emph{śobhate} is the present of the root √śubh, \textquotedblleft{}to beautify; to look beautiful, shine\textquotedblright{} (MW). The adjective \emph{śubha}, \textquotedblleft{}bright, auspicious\textquotedblright{}, comes from the same root, as in Hindi \emph{śubh rātri}, \textquotedblleft{}good night\textquotedblright{}.
\end{cousinweb}

\begin{grammarlens}[title={What you've built}]
One more word for this chapter.
\end{grammarlens}

\subsection*{Your turn}
\begin{itemize}
\raggedright
  \item \emph{Say it:} \emph{śobhate}
  \item \emph{Say it:} \emph{śobhate}, once more
  \item \emph{Say it:} say \emph{śobhate}
  \item \emph{Recall:} say the Sanskrit for begs, then the Sanskrit for speaks, then say \emph{śobhate} again
\end{itemize}

\begin{tcolorbox}[breakable,colback=teal!4,colframe=teal!35!black,title={Before you move on}]
What does \sk{शोभते} mean? (\textquotedblleft{}Shines, looks fine\textquotedblright{}.) Say it once more.
\end{tcolorbox}
